\documentclass[12pt]{article}
\usepackage[margin=0.6in,top=0.65in,bottom=0.5in]{geometry}
\usepackage{lmodern}
\usepackage{microtype}
\usepackage{fancyhdr}
\usepackage{titlesec}
\usepackage{enumitem}
\usepackage{lastpage}
\usepackage[hidelinks]{hyperref}

\setlength{\parindent}{0pt}
\setlength{\parskip}{0.22em}
\setlength{\headheight}{26pt}
\raggedbottom
\titleformat{\section}{\normalsize\bfseries\scshape}{}{0pt}{}
\titleformat{\subsection}{\normalsize\bfseries}{}{0pt}{}
\titlespacing*{\section}{0pt}{0.4em}{0.15em}
\titlespacing*{\subsection}{0pt}{0.3em}{0.1em}
\pagestyle{fancy}
\fancyhf{}
\fancyhead[C]{\footnotesize Lit Example Reader \quad Assignment ID: U1L6BLE \quad Page \thepage\ of \pageref{LastPage}}
\renewcommand{\headrulewidth}{0.5pt}

\newenvironment{playpassage}{\begin{quote}\small\setlength{\parskip}{0.12em}\setlength{\parindent}{0pt}}{\end{quote}}
\newcommand{\speaker}[1]{\textbf{#1.}}

\begin{document}

\begin{center}
{\LARGE\bfseries Week 6B Lit Example Reader}\par\vspace{0.02em}
{\large\scshape Either/Or and If/Then in Caesar and Macbeth}\par\vspace{0.08em}
\rule{0.78\textwidth}{0.5pt}
\end{center}

\textbf{Purpose:} Literary passages for Week 6B: test compound pressure---either/or alternatives and if/then conditions---without treating persuasion as proof.

\textbf{What to watch for:} What options are offered? What condition is assumed? What follows if the compound claim is granted? What remains unproved?

\section*{Before the Passages: The Week 6B Logic Tools}

\begin{itemize}[leftmargin=1.3em,itemsep=0.08em,topsep=0.1em,parsep=0.04em]
\item \textbf{Either/or:} alternatives offered as if they settle a choice.
\item \textbf{False choice:} an either/or that ignores other live options.
\item \textbf{If/then:} a condition linked to a result (antecedent $\rightarrow$ consequent).
\item \textbf{Valid moves:} grant the if-part to get the then-part; deny the then-part to reject the if-part.
\item \textbf{Traps:} treating the then-part as proof of the if-part; treating denial of the if-part as proof that the then-part fails.
\item \textbf{Form vs fact:} even a clean compound pattern fails if a premise is false, loaded, or unproved.
\end{itemize}

\section*{Passage 1: Excerpts from \textit{Julius Caesar}, Act 3, Scene 2 --- Brutus's Either/Or}

\textbf{Context:} After Caesar's death, Brutus addresses the Roman crowd. Watch how he frames the choice between Caesar living and the people living free.

\begin{playpassage}
\speaker{Brutus} If there be any in this assembly, any dear friend of Caesar's, to him I say that Brutus' love to Caesar was no less than his. If then that friend demand why Brutus rose against Caesar, this is my answer:---Not that I loved Caesar less, but that I loved Rome more. Had you rather Caesar were living and die all slaves, than that Caesar were dead, to live all free men? As Caesar loved me, I weep for him; as he was fortunate, I rejoice at it; as he was valiant, I honour him: but, as he was ambitious, I slew him. There is tears for his love; joy for his fortune; honour for his valour; and death for his ambition.
\end{playpassage}

\subsection*{Reader Notes}

Here is your most assisted example. Brutus pairs love of Caesar with love of Rome, then presses a hard alternative: Caesar living with the people as slaves, or Caesar dead with the people free. The speech can feel finished because the either/or is vivid and moral. A disciplined listener still separates several jobs: (1) what either/or is being offered; (2) whether those two options really exhaust the live possibilities; (3) whether words such as \textit{slave}, \textit{free}, and \textit{ambitious} have been fixed; and (4) whether ambition has been proved, not merely asserted. Compound form can create pressure even before the fact test is passed.

\textbf{Reading task:} Underline the either/or question. Circle \textit{slaves}, \textit{free}, and \textit{ambitious}. In the margin, write one other live option a careful Roman might still want considered.

\newpage

\section*{Passage 2: Excerpts from \textit{Macbeth}, Act 1, Scenes 3 and 5 --- Conditional Prophecy Pressure}

\textbf{Context:} The weird sisters greet Macbeth with titles he has not yet fully secured. Later Lady Macbeth receives news of the prophecy. Watch if/then pressure: what is treated as condition, and what is treated as already proved.

\begin{playpassage}
\speaker{First Witch} All hail, Macbeth! hail to thee, thane of Glamis!\\
\speaker{Second Witch} All hail, Macbeth, hail to thee, thane of Cawdor!\\
\speaker{Third Witch} All hail, Macbeth, thou shalt be king hereafter!\\[0.35em]
\speaker{Macbeth} \textit{(aside)} Glamis, and thane of Cawdor!\\
The greatest is behind.\\[0.35em]
\speaker{Lady Macbeth} Glamis thou art, and Cawdor; and shalt be\\
What thou art promised: yet do I fear thy nature;\\
It is too full o' the milk of human kindness\\
To catch the nearest way: thou wouldst be great;\\
Art not without ambition, but without\\
The illness should attend it\ldots\\
Hie thee hither,\\
That I may pour my spirits in thine ear;\\
And chastise with the valour of my tongue\\
All that impedes thee from the golden round.
\end{playpassage}

\subsection*{Reader Notes}

This passage shifts more work to you. Notice how later certainty grows from earlier greeting and promise. A listener can reconstruct pressures such as: if the sisters speak truly, then kingship is ahead; if Cawdor has come true, then the greater promise is nearer; if the crown is promised, then hesitation is waste. Mark where a condition is treated as if it were already proof. Decide whether a fulfilled smaller title proves the larger destiny, or only makes a conditional hope feel stronger.

\textbf{Reading task:} Star any line that functions like an if/then. In the margin, note one place where a later fact might be used as if it proved an earlier prophecy. Do not yet write the full worksheet diagnosis; gather evidence only.

\end{document}
